\chapter{PHÂN TÍCH VÀ THIẾT KẾ HỆ THỐNG}

\section{Phân tích yêu cầu hệ thống}
\subsection{Yêu cầu chức năng (Functional Requirements)}
% Viết phân tích chi tiết các nhóm chức năng tại đây (Khoảng 600 từ)
\begin{itemize}
    \item Nhóm chức năng kiểm soát và nhận diện tại thiết bị biên.
    \item Nhóm chức năng cấu hình và quản trị thông qua Mobile App.
    \item Nhóm chức năng an ninh, giám sát và cảnh báo hệ thống.
\end{itemize}

\subsection{Yêu cầu phi chức năng (Non-functional Requirements)}
% Viết đặc tả chỉ số chất lượng phần cứng, phần mềm (Khoảng 400 từ)
% Chèn bảng đặc tả hiệu năng ở đây
\begin{table}[htbp]
\caption{Đặc tả các chỉ số phi chức năng của hệ thống}
\label{tab:non_functional}
\centering
\begin{tabular}{|l|l|l|}
\hline
\textbf{STT} & \textbf{Tiêu chí chất lượng} & \textbf{Ngưỡng thông số kỹ thuật} \\ \hline
1 & Thời gian phản hồi nhận diện & $< 1.5$ giây \\ \hline
2 & Tỷ lệ chấp nhận sai người (FAR) & $< 0.001\%$ \\ \hline
3 & Tỷ lệ từ chối đúng người (FRR) & $< 1\%$ \\ \hline
4 & Tính sẵn sàng của hệ thống (Availability) & Chạy độc lập 24/7 kể cả mất Internet \\ \hline
\end{tabular}
\end{table}

\section{Sơ đồ kiến trúc hệ thống (System Architecture)}
\subsection{Kiến trúc tổng thể hệ thống mạng Edge-to-Cloud}
% Phân tích mối quan hệ giữa ESP32-S3, Server NestJS và App Flutter (Khoảng 500 từ)
% \begin{figure}[htbp]
% \centering
% \includegraphics[width=0.8\textwidth]{images/ch3/system_architecture.png}
% \caption{Sơ đồ kiến trúc tổng thể hệ thống}
% \end{figure}

\subsection{Sơ đồ khối chức năng phần cứng tại biên}
% Đặc tả luồng tương tác linh kiện ngoại vi kết nối với ESP32-S3 (Khoảng 500 từ)
% \begin{figure}[htbp]
% \centering
% \includegraphics[width=0.8\textwidth]{images/ch3/hardware_block.png}
% \caption{Sơ đồ khối chức năng phần cứng thiết bị khóa}
% \end{figure}

\section{Các biểu đồ mô hình hóa phần mềm UML}
\subsection{Biểu đồ Use Case tổng thể và đặc tả Use Case}
% Mô tả các Actor và các ca sử dụng tương ứng (Khoảng 600 từ)
% \begin{figure}[htbp]
% \centering
% \includegraphics[width=0.7\textwidth]{images/ch3/usecase_diagram.png}
% \caption{Biểu đồ Use Case tổng thể của hệ thống}
% \end{figure}

\subsection{Biểu đồ Hoạt động (Activity Diagram)}
% Giải thích luồng xử lý thuật toán nhận dạng rẽ nhánh điều kiện (Khoảng 900 từ)
% \begin{figure}[htbp]
% \centering
% \includegraphics[width=0.6\textwidth]{images/ch3/activity_recognition.png}
% \caption{Biểu đồ hoạt động luồng xử lý nhận diện khuôn mặt mở cửa}
% \end{figure}

\subsection{Biểu đồ Trình tự (Sequence Diagram)}
% Mô tả thông điệp truyền tải thời gian thực đa giao thức HTTP/WebSocket/MQTT (Khoảng 900 từ)
% \begin{figure}[htbp]
% \centering
% \includegraphics[width=0.9\textwidth]{images/ch3/sequence_remote_open.png}
% \caption{Biểu đồ trình tự chức năng ra lệnh mở cửa từ xa từ Mobile App}
% \end{figure}

\section{Thiết kế cơ sở dữ liệu (Database Design)}
\subsection{Thiết kế cấu trúc lưu trữ nhúng cục bộ trên bộ nhớ Flash}
% Đặc tả phân vùng NVS, tính toán kích thước ô nhớ 512 dimensions float vector (Khoảng 400 từ)

\subsection{Thiết kế mô hình dữ liệu quan hệ tập trung trên Server}
% Phân tích cấu trúc các bảng thực thể (Users, Tenants, Occupants, AccessLogs...) 
% Giải thích logic ánh xạ mối quan hệ thừa kế (R8a) và quan hệ 1-N (R7) (Khoảng 600 từ)
% \begin{figure}[htbp]
% \centering
% \includegraphics[width=1.0\textwidth]{images/ch3/erd_diagram.png}
% \caption{Sơ đồ quan hệ thực thể ERD hệ thống quản lý tập trung}
% \end{figure}